%! Inverse Functions — Unit 1, Lesson 1.6 — Homework (Answer Key)
\documentclass[10pt]{article}
\usepackage{precalculus-article}
\usepackage{precalculus-key}   % pulls in -boxes; do NOT also load -boxes
\newcommand{\TallMath}[1]{$\displaystyle #1\rule[-1.4em]{0pt}{3.2em}$}

\begin{document}

\pageheader{Unit 1, Lesson 1.6}{Homework --- Answer Key}

\vspace{0.12in}

\begin{enumerate}[itemsep=10pt]

\item Two functions are given by tables.

\smallskip
\renewcommand{\arraystretch}{1.4}
\begin{center}
\begin{tabular}{|c|c|c|c|c|}
\hline
$x$ & 0 & 1 & 2 & 3 \\
\hline
$f(x)$ & 7 & 4 & 1 & $-2$ \\
\hline
\end{tabular}
\qquad
\begin{tabular}{|c|c|c|c|c|}
\hline
$x$ & 0 & 1 & 2 & 3 \\
\hline
$g(x)$ & 5 & 2 & 2 & 9 \\
\hline
\end{tabular}
\end{center}

\begin{enumerate}[label=(\alph*), itemsep=4pt]
  \item Complete the table for $f^{-1}$ by trading rows.

        \smallskip
        \renewcommand{\arraystretch}{1.4}
        \begin{center}
        \begin{tabular}{|c|c|c|c|c|}
        \hline
        $x$ & 7 & 4 & 1 & $-2$ \\
        \hline
        $f^{-1}(x)$ & \ans{0} & \ans{1} & \ans{2}
          & \ans{3} \\
        \hline
        \end{tabular}
        \end{center}

  \item Is $g$ invertible? \ans{No}

        Why or why not? \ans{$g(1) = g(2) = 2$, so one output has two inputs.}
\end{enumerate}

\item Find the inverse of each function by swapping and solving.

\begin{work}
              y &= 5x + 3 \\
              x &= 5y + 3 \\
          x - 3 &= 5y \\
  \frac{x-3}{5} &= y \\
              y &= \frac{x}{2} - 7 \\
              x &= \frac{y}{2} - 7 \\
          x + 7 &= \frac{y}{2} \\
       2(x + 7) &= y
\end{work}

\begin{enumerate}[label=(\alph*), itemsep=4pt]
  \item $f(x) = 5x + 3$ \qquad $f^{-1}(x) = $ \ans{$(x-3)/5$}
  \item $g(x) = \dfrac{x}{2} - 7$ \qquad $g^{-1}(x) = $ \ans{$2(x+7)$}
\end{enumerate}

\item Use your answer to 2(a) to verify the pair \textbf{in both directions}.

\begin{work}
  f\big(f^{-1}(x)\big) &= 5\left(\frac{x-3}{5}\right) + 3 \\
                       &= (x - 3) + 3 \\
                       &= x \\
  f^{-1}\big(f(x)\big) &= \frac{(5x + 3) - 3}{5} \\
                       &= \frac{5x}{5} \\
                       &= x
\end{work}

$f\big(f^{-1}(x)\big) = $ \ans{$x$} \qquad
$f^{-1}\big(f(x)\big) = $ \ans{$x$}

\boxguard[20]
\item \textbf{Back to the thermometer.} The lab's conversion rule is
      $F(x) = 1.8x + 32$, taking degrees Celsius to degrees Fahrenheit.

\begin{work}
                     y &= 1.8x + 32 \\
                     x &= 1.8y + 32 \\
                x - 32 &= 1.8y \\
     \frac{x - 32}{1.8} &= y \\
           F^{-1}(212) &= \frac{212 - 32}{1.8} \\
                       &= 100
\end{work}

\begin{enumerate}[label=(\alph*), itemsep=4pt]
  \item $F^{-1}(x) = $ \ans{$(x-32)/1.8$}
  \item $F^{-1}(212) = $ \ans{100}. In one sentence, what does that number
        mean? \ans{Water boils at 100 degrees Celsius.}
  \item $F(50) = $ \ans{122}, so $\dfrac{1}{F(50)} = $ \ans{$1/122$},
        while $F^{-1}(50) = $ \ans{10}. Circle one:
        $F^{-1}(50)$ \textbf{is} / \ans{is not} the same as
        $\dfrac{1}{F(50)}$.
\end{enumerate}

\item $f(x) = x^2$ has no inverse on all real numbers.

\begin{enumerate}[label=(\alph*), itemsep=4pt]
  \item Name two different inputs with the same output:
        \ans{$-3$} and \ans{3} both give \ans{9}.
  \item Restrict the domain to \ans{$x \ge 0$} so that $f$ becomes one-to-one.
  \item On that restricted domain, $f^{-1}(x) = $ \ans{$\sqrt{x}$}, with domain
        \ans{$x \ge 0$}.
\end{enumerate}

\item A function $f$ satisfies $f(3) = -2$. Which statement must be true?

\begin{enumerate}[label=\Alph*., itemsep=3pt]
  \item $f^{-1}(3) = -2$
  \item $f^{-1}(-2) = 3$ \quad \textcolor{keyred}{\textbf{$\leftarrow$ correct}}
  \item $f^{-1}(3) = -\dfrac{1}{2}$
  \item $f^{-1}(-2) = -\dfrac{1}{3}$
\end{enumerate}

\end{enumerate}

\vspace{0.1in}

\boxguard
\begin{extensionbox}
Let $f(x) = x^2$ on all real numbers and let $h(x) = \sqrt{x}$. Compute
$f\big(h(9)\big)$ and $h\big(f(-3)\big)$. One of them returns the number you
started with and the other does not. Explain which composition fails and what
domain restriction on $f$ repairs the claim that $f$ and $h$ are inverses.

\ansline{$f(h(9)) = (\sqrt{9})^2 = 9$, but $h(f(-3)) = \sqrt{9} = 3$, not $-3$.}
\ansline{The second fails because $f$ is not one-to-one; restricting $f$ to $x \ge 0$ repairs it.}
\end{extensionbox}

\vspace{0.1in}

\boxguard
\begin{notesbox}{Preview of Lesson 1.7}
Today you learned to run a function backwards. Lesson 1.7, \textit{Function
Model Selection and Construction}, asks the prior question: which kind of
function should model a situation in the first place? Once a model is chosen,
inverting it is how you answer ``what input produced this output?'' Much later,
in Unit 5, the logarithm arrives as nothing more than the inverse of the
exponential function --- built with exactly today's method.
\end{notesbox}

\end{document}
